\documentclass[11pt,letterpaper]{article}
\usepackage[T1]{fontenc}
\usepackage[utf8]{inputenc}
\usepackage{lmodern}
\usepackage[margin=1in,headheight=15pt]{geometry}
\usepackage{microtype,amsmath,amssymb,amsthm,mathtools}
\usepackage{booktabs,longtable,array,tabularx}
\usepackage{enumitem,listings,xcolor,xurl}
\usepackage{fancyhdr,hyperref}
\hypersetup{colorlinks=true,linkcolor=black,citecolor=black,urlcolor=blue,
 pdftitle={Forge: Three Algorithmic Extensions Beyond the Current Design},
 pdfauthor={OpenAI},pdfsubject={Invariant subspaces, telescoping, and noncommutative certificates}}
\pagestyle{fancy}\fancyhf{}
\fancyhead[L]{\small Forge: algorithmic extensions}
\fancyhead[R]{\small September 15, 2026}
\fancyfoot[C]{\thepage}
\setcounter{tocdepth}{1}
\setlength{\parskip}{4pt}\setlength{\parindent}{0pt}
\setlist{itemsep=2pt,topsep=4pt,leftmargin=1.5em}
\lstset{basicstyle=\ttfamily\small,columns=fullflexible,keepspaces=true,
 breaklines=true,frame=single,framerule=.3pt,showstringspaces=false,
 xleftmargin=3pt,xrightmargin=3pt,aboveskip=8pt,belowskip=8pt}
\newcommand{\Q}{\mathbb Q}\newcommand{\R}{\mathbb R}\newcommand{\N}{\mathbb N}
\newcommand{\Z}{\mathbb Z}\newcommand{\Span}{\operatorname{span}}
\newcommand{\rank}{\operatorname{rank}}\newcommand{\im}{\operatorname{im}}
\newcommand{\forge}{\texttt{forge}}\newcommand{\grind}{\texttt{grind}}
\newcommand{\code}[1]{\texttt{#1}}
\newcommand{\unknown}{\textsc{unknown}}
\newcommand{\status}[1]{\par\smallskip\noindent\textbf{Evidence status.} #1\par\smallskip}
\newtheorem{theorem}{Theorem}[section]
\newtheorem{proposition}[theorem]{Proposition}
\newtheorem{lemma}[theorem]{Lemma}
\newtheorem{corollary}[theorem]{Corollary}
\theoremstyle{definition}\newtheorem{definition}[theorem]{Definition}
\newtheorem{example}[theorem]{Example}
\theoremstyle{remark}\newtheorem{remark}[theorem]{Remark}
\DeclareMathOperator{\degw}{deg}
\begin{document}
\begin{titlepage}
\vspace*{0.6in}
{\Large\sffamily RESEARCH AND IMPLEMENTATION REPORT\par}
\vspace{0.35in}
{\Huge\bfseries Forge:\\[6pt] Three Algorithmic Extensions\\[6pt] Beyond the Current Design\par}
\vspace{0.3in}
{\Large Inductive subspaces, boundary-aware telescoping,\\and two-sided noncommutative certificates\par}
\vspace{0.45in}
{\large Prepared for Vladimir Reshetnikov\par}
{\large September 15, 2026\par}
\vspace{0.35in}
\rule{\textwidth}{0.5pt}
\vspace{0.1in}
\textbf{The practical recommendation.} Keep Forge's existing orchestration design.
Add three narrow workers that discover mathematical structure the current
merged proposal does not supply algorithmically: invariant spaces closed under
transitions, hypergeometric antidifferences and telescopers, and consequences of
noncommutative relations with multiplication on both sides.

\textbf{Delivered evidence.} Executed Python search and independent exact replay
for 65 stored certificates; 188 rejected certificate mutations; 1,100 assertions
in the generated experiment; 13 named regression tests; reproducible source,
problem data, certificates, and logs. The invariant ablation solves 24 of 24
selected cases, versus 3 of 24 for the same-feature conservation-only reference.

\textbf{Verification boundary.} These are research prototypes, not a completed
Lean tactic. No Lean compiler was available in this execution environment.
Included Lean files are uncompiled replay targets, not kernel-checked results.
In particular, polynomial telescoper replay is not silently promoted to a proof
of a finite sum with singular endpoints.
\vfill
{\small Source baseline: Forge commit
\nolinkurl{674521027d968d59f7b83220ed52304a85cb55e2}.\\
The additions are relative to the inspected design, not claims of new invention
of invariant theory, Gosper's algorithm, or noncommutative ideal theory.}
\end{titlepage}

\begin{abstract}
The existing Forge repository already contains an unusually broad architecture:
obligation-directed search, induction planning, exact arithmetic certificates,
witness synthesis, selective equality exploration, and carefully delimited
external search. Repeating that architecture would not make the tactic more
capable. This report develops three additional, executable mathematical workers.
First, a descending sequence of rational linear subspaces discovers families of
polynomial invariants satisfying $p(T(x))=H p(x)$, allowing scaling and mutual
dependence rather than requiring each polynomial to be conserved. It terminates
by dimension descent and is complete for bounded-degree polynomial equalities
of the specified affine, nondeterministically branching fragment. Second,
coefficient-linear searches discover hypergeometric antidifferences and an
order-one creative telescoper. A worked squared-binomial example makes the
singular-boundary problem explicit and derives a nonsingular flux, so that
algebraic certificate checking and finite-sum semantics remain distinct.
Third, bounded two-sided relation search in a free associative algebra emits
identities $f=\sum c\,u r v$, enabling matrix and operator reasoning without
assuming commutativity or finding a terminating rewrite orientation.
Each worker comes with a concrete search algorithm, a separate exact checker,
mathematical soundness arguments, sharply bounded completeness claims, and a
Lean integration route. The accompanying experiments include positive and
negative controls, adversarial mutations, an invariant ablation, and
standard-library-only replay in a fresh Python process. The report also explains
how these workers can strengthen universal behavioral checking of exact
Leant/Djex candidates without redesigning their existing typed streams.
\end{abstract}
\tableofcontents
\newpage

\section{Recommendation and scope}
\label{sec:recommendation}

\subsection{Add mathematical discovery, not another controller}
Forge's next useful increment should be a small number of workers with a clear
answer to four questions: what do they discover, what finite object do they
return, what does its checker establish, and how does the resulting theorem
help a blocked goal? The three proposed workers answer those questions without
replacing the existing obligation graph or inventing a second policy for proof
admission.

The invariant worker discovers \emph{relations preserved collectively}. The
summation worker discovers \emph{an antidifference or a recurrence with a local
flux}. The noncommutative worker discovers \emph{which left and right contexts
make the given relations add up to the target}. These are different ways of
finding an intermediate proof object, not merely running a familiar tactic with
more fuel.

The intended deployment is initially conservative. A worker receives a supported
problem extracted from the actual Lean goal. Search proposes a certificate.
Ordinary Lean theorem applications and existing normalizers reconstruct its
small algebraic identities. The new local theorem is then given to Forge's
existing machinery. A reflected checker can replace large proof scripts later;
there is no need to wait for a new kernel-facing plugin architecture before
measuring the first useful cases.

\subsection{What was actually inspected}
The review is pinned to Forge commit
\nolinkurl{674521027d968d59f7b83220ed52304a85cb55e2}, Leant commit
\nolinkurl{6bf05ad78c467989e68290f2d08bbed40802d485}, and Djex commit
\nolinkurl{e8778f4ebd63e1f9b9b410fa4de8d14a8a04c9e5}.
The Forge README, repository tree, merged structural, nonlinear, Boolean,
integration, and provenance sections, and its Leant/Djex review supplied the
baseline. The neighboring projects' behavioral-synthesis and streaming reports
were also read at the listed revisions~\cite{forge,forge-structural,forge-nonlinear,
forge-external,forge-integration,forge-provenance,leant-behavior,djex-stream}.

A complete local clone was not available. Targeted repository searches for
summation-related and noncommutative terms returned no matches, but search
absence is not a proof that no historical file contains an adjacent idea.
Consequently the nonduplication claim is deliberately scoped: the algorithms
below extend the \emph{inspected merged design and its provenance ledger}.
This is not a line-by-line review of every original proposal or of all Haskell
implementation files.

At the reviewed revision, Forge describes a design and a collection of
prototypes, not an implemented all-purpose Lean tactic. The repository's own
Lean experiments are distinct from the experiments delivered here. Its baseline
identifies Lean v4.34.0 and Mathlib commit
\nolinkurl{1cf325a0cf67aca2b04d76b5380ff6a9e410aefa}; this report reuses those as
\emph{integration targets}, not as locally tested dependencies~\cite{forge,
forge-integration}.

\subsection{A nonduplication ledger}
\begin{table}[tbp]\small\centering
\begin{tabular}{@{}p{.27\textwidth}p{.31\textwidth}p{.35\textwidth}@{}}
\toprule
Area & Already in the inspected design & Increment developed here\\
\midrule
Inductive equations & Polynomial finite-difference CEGIS; conserved state
quantities; a warning that jointly unknown invariant ideals lead to bilinear
search & A dimension-descending subspace algorithm, transition matrices,
finite-fragment completeness, and a same-feature ablation\\
Finite sums & Additive polynomial recurrences and exact coefficient checking &
Rational hypergeometric multipliers; an explicitly searched bivariate
telescoper; nonsingular flux and boundary obligations\\
Algebra & Commutative polynomial cones, equality ideals, and selective
rewriting/e-graphs & Two-sided consequences in a free associative algebra, word
order preserved, with coefficient-linear search and relation witnesses\\
Leant/Djex transfer & Typed streams, exact-candidate verification, evidence
ownership, fairness, and negative-evidence restrictions & A specific universal
behavioral-proof backend and product-state extraction contract for eligible
candidates\\
\bottomrule
\end{tabular}
\end{table}

The controller, Horn slicing, induction-variable selection, accumulator
generalization, anti-unification, SOS/Bernstein/Sturm machinery, Farkas and
lattice witnesses, CDCL(T), basic LJT proposals, and generic security boundaries
are not reintroduced as contributions. They are infrastructure or neighboring
workers that these additions can use~\cite{forge-structural,forge-nonlinear,
forge-external,forge-integration}.

\subsection{Four evidence levels used throughout}
A \emph{mathematical theorem} below has an ordinary written proof. An
\emph{executed algorithm} has a runnable implementation and a recorded result.
An \emph{exact replay} has been checked by the supplied Python checker against
the separately supplied source problem. A \emph{Lean proof} would additionally
require successful elaboration and kernel checking at the actual pin.
Only the first three levels were attained in this run. Source code ending with
\code{by ring} is not reported as a fourth-level result merely because it looks
plausible.

This distinction is especially important for telescoping: the current checker
accepts an identity in a rational polynomial ring. The theorem connecting that
identity to a source sequence, its index domain, and its boundary conditions
must still be proved. The experiment does not disguise this gap with a
misleading universal \textsc{proved} status.

\section{Inductive subspaces instead of conserved scalars}
\label{sec:invariant}
\status{General finite-feature search and exact closure replay are implemented.
Twenty-four positive cases, three expected-unknown cases, a same-feature
conservation ablation, and corruption tests were executed. Lean extraction and
a generic reflected checker are not implemented.}

\subsection{The capability gap}
A conserved quantity satisfies $p(T(x))=p(x)$. An inductive relation needs much
less: whenever $p(x)=0$, its updated value must also be zero. For example,
\[
 T(x,y,z)=(2x,3y,6z),\qquad p(x,y,z)=z-xy
\]
satisfies $p(T(x,y,z))=6p(x,y,z)$. Starting at $(u,v,uv)$, the relation
$z=xy$ remains true after every number of steps. A search restricted to
$p\circ T=p$ throws this relation away for a reason unrelated to its validity.

There is a second gap: a useful family can be preserved while none of its
chosen basis elements is preserved in isolation. Under $T(x,y)=(y,2x)$, the
family $(y,x)$ transforms by
\[
 \begin{pmatrix}y\circ T\\x\circ T\end{pmatrix}
 =\begin{pmatrix}0&2\\1&0\end{pmatrix}
  \begin{pmatrix}y\\x\end{pmatrix}.
\]
Both vanish forever from $(0,0)$, although a scalar conservation test finds no
nonzero linear quantity of that kind.

Forge's structural section already recognizes the general invariant problem
and warns that simultaneously solving for $p_i$ and arbitrary polynomial
multipliers in $p_i\circ T=\sum_j h_{ij}(x)p_j$ is bilinear. That warning is
correct. The algorithm here does not rename that bilinear problem ``linear.''
It chooses a different, useful fragment and computes its greatest fixed
subspace by successive linear problems~\cite{forge-structural}.

\subsection{Problem model}
Let $K$ be an infinite field of characteristic zero; the prototype calculates
coefficients in $\Q$. State variables are $x=(x_1,\ldots,x_n)$. Initial states
are specified by a polynomial map
\[
 I:K^p\longrightarrow K^n,
\]
including $p=0$ for a single fixed state. A finite nonempty collection of
polynomial maps $T_b:K^n\to K^n$ gives the permitted updates. Every branch is
initially treated as nondeterministic: it may be taken at any reachable state.

Reachable states are obtained from some $I(u)$ by a finite composition of the
$T_b$. A guard can be ignored to obtain a sound overapproximation; it cannot be
silently used as a proof assumption. A guard-sensitive refinement is discussed
later. There is no assertion here that the transition system terminates.

Choose a finite-dimensional rational space
\[
 E=\Span_{\Q}\{m_1,\ldots,m_N\}\subseteq\Q[x].
\]
The implementation uses all monomials of total degree at most $D$, so
$N=\binom{n+D}{D}$. The mathematical construction also works for a selected
feature list after removing linear dependencies. The pullback
$T_b^*p=p\circ T_b$ is linear in $p$, even when $T_b$ is nonlinear.
This observation supplies the useful linear structure.

\subsection{Descending closure}
Define
\begin{align}
 V_0&=\{p\in E:p\circ I=0\text{ as a polynomial}\},\label{eq:v0}\\
 V_{j+1}&=V_j\cap\bigcap_b(T_b^*)^{-1}(V_j).
 \label{eq:vnext}
\end{align}
The inverse-image condition means that $p\circ T_b$ belongs to the current
subspace as a polynomial, not merely after discarding high-degree monomials.
At each iteration the current basis is fixed; only the coefficients of a new
candidate and its coordinates in that fixed basis are unknown. These are linear
constraints.

\begin{theorem}[Termination and maximality]
\label{thm:subspace}
The sequence \eqref{eq:vnext} stabilizes after at most $N$ strict dimension
drops. Its final space $V_*$ is the largest subspace of $E$ that vanishes under
initialization and is stable under every pullback $T_b^*$.
Every polynomial in $V_*$ vanishes at every reachable state.
\end{theorem}
\begin{proof}
Each $V_j$ is a linear subspace, and $V_{j+1}\subseteq V_j$. A strict inclusion
between finite-dimensional spaces strictly decreases dimension, so there can
be at most $N$ such inclusions. Equal dimensions at a step imply equality of
spaces and therefore a fixed point.

Suppose $W\subseteq E$ vanishes under initialization and $T_b^*W\subseteq W$
for every branch. Then $W\subseteq V_0$. If $W\subseteq V_j$, every $p\in W$
also has $T_b^*p\in W\subseteq V_j$, whence $p\in V_{j+1}$. Induction gives
$W\subseteq V_*$, proving maximality.

At an initial state every $p\in V_*$ vanishes. At a transition from a state
where all members of $V_*$ vanish, each $p\circ T_b$ is a linear combination of
members of $V_*$. It therefore vanishes too. Induction on the finite execution
length proves the last assertion.
\end{proof}

The conclusion does not require the returned basis to be linearly independent
for \emph{soundness}. Independence is a search and size property; a checker
can accept a redundant family if all displayed identities are valid.

\subsection{Exact matrix construction}
The coefficient representation must include both the original features and all
their pullbacks. Let $\mathcal M$ be the union of monomial supports of $E$ and
$T_b^*E$ for every branch, and let $M=|\mathcal M|$. In this common ambient
coordinate space, write
\[
 A_E\in\Q^{M\times N},\qquad A_b\in\Q^{M\times N}
\]
for the columns representing $m_i$ and $m_i\circ T_b$ respectively. Write
$C_I$ for the coefficient matrix of the $m_i\circ I$. Then a basis matrix for
$V_0$ is a nullspace basis of $C_I$.

Suppose $C_j\in\Q^{N\times d_j}$ has columns spanning $V_j$. Compute a matrix
$L_j$ whose rows span the left nullspace of $A_EC_j$. Membership in that column
space is equivalent to annihilation by $L_j$. Therefore
\begin{equation}
 C_{j+1}=C_j\,\operatorname{basisKer}
 \begin{pmatrix}
  L_jA_1C_j\\ \vdots\\ L_jA_BC_j
 \end{pmatrix}.
 \label{eq:matrixstep}
\end{equation}
All matrices in a single call to \code{basisKer} are already known. There are
no products of two independently unknown coefficient families.

\par\noindent\begin{minipage}{\linewidth}
\begin{lstlisting}
features := monomials(state_variables, degree_cap)
C := kernel(coefficients(features after initialization))
ambient := supports(features and every branch pullback)
AE, Ab := coefficient matrices in that same ambient
repeat:
    L := transpose(kernel(transpose(AE * C)))
    K := kernel(stack_b(L * Ab * C))
    Cnext := independent_columns(C * K)
    if columns(Cnext) == columns(C): return Cnext
    C := Cnext
\end{lstlisting}
\end{minipage}\par\smallskip

The ambient support is \emph{not} recomputed by truncating pullbacks to degree
$D$. Such truncation would make $x\mapsto x^2$ look like the zero map on the
space spanned by $x$ and would destroy the meaning of the test.
Simultaneous substitution is equally important: replacing $(x,y)$ by $(y,x)$
must not become two successive substitutions that accidentally identify both
variables. The implementation and regression suite test this explicitly.

\subsection{Certificate and independent checker}
If the final basis is $p=(p_1,\ldots,p_r)^\mathsf T$, search returns a matrix
$H_b\in\Q^{r\times r}$ per branch and a vector $a\in\Q^r$ for a requested
target $g$, with
\begin{align}
 p_i(I(u))&=0,\label{eq:initcert}\\
 p(T_b(x))&=H_bp(x),\label{eq:stepcert}\\
 g(x)&=a^\mathsf T p(x).\label{eq:targetcert}
\end{align}
The checker obtains $I,T_b,g$ from the separately provided problem. It expands
\eqref{eq:initcert}--\eqref{eq:targetcert} exactly. It neither reruns the
fixed-point algorithm nor accepts a target polynomial embedded by the producer
as a replacement for the source target.

\begin{corollary}[Certificate soundness]
An accepted certificate proves $g(s)=0$ at every reachable state of the supplied
transition system.
\end{corollary}
\begin{proof}
The initialization and transition identities prove that all components of $p$
vanish by induction. Equation~\eqref{eq:targetcert} gives the target.
\end{proof}

The prototype producer uses SymPy expressions, polynomial coefficient matrices,
and exact nullspaces. The checker instead uses sparse exponent tuples and
Python's standard-library \code{Fraction}. This separates matrix construction
and solving from checking the result. Both programs still depend on the Python
runtime and rational arithmetic; algorithmic independence is not a formal
verification theorem.

\subsection{A useful completeness theorem, and its exact boundary}
\begin{theorem}[Bounded-degree affine completeness]
\label{thm:affinecomplete}
Suppose every $T_b$ is affine, $E$ contains all monomials of total degree at most
$D$, and initial states are exactly the image of the supplied polynomial map
$I:K^p\to K^n$. Ignoring resource limits, $V_*$ consists of precisely the
rational polynomials of degree at most $D$ that vanish at every reachable
state.
\end{theorem}
\begin{proof}
An affine substitution maps a polynomial of degree at most $D$ to one of degree
at most $D$, so every pullback preserves $E$. Let $W$ be the space of members
of $E$ vanishing on the reachable set. It vanishes on every initial state.
Because $K$ is infinite, a polynomial in the initial parameters vanishing at all
parameter values is the zero polynomial; hence $W\subseteq V_0$.
If $p\in W$, then $p(T_b(s))=0$ for every reachable $s$, since its successor is
reachable. Thus $T_b^*p\in W$. Maximality gives $W\subseteq V_*$. Soundness
gives the reverse inclusion.
\end{proof}

This is not the strongest known result about affine programs. Algorithms for
strongest algebraic invariants, without this fixed degree-space restriction,
are established prior work~\cite{affine-invariants}. The merit here is the
small, executable fragment and its cheap identity certificate, not a claim to
have solved invariant generation for the first time.

For a selected feature set not closed under pullback, the maximality theorem
still holds but affine completeness may not. For nonlinear transitions the
failure can be unavoidable at every finite degree.

\begin{example}[A nonlinear invariant that no finite closed subspace sees]
Let $T(x)=x^2$, initialize $x=0$, and take $E=\Q[x]_{\le D}$ for any finite
positive $D$. The reachable state is always zero, so $x=0$ is an invariant.
Nevertheless, any nonzero pullback-stable subspace vanishing at zero would
contain a polynomial of maximal positive degree $d$. Its pullback has degree
$2d>d$, a contradiction. Hence $V_*=0$ for every finite $D$.
The ordinary ideal certificate $x\circ T=x\cdot x$ succeeds, but its multiplier
is polynomial rather than constant.
\end{example}

The right response is to route this case to a candidate-fixed multiplier
search or a different induction argument, not to advertise arbitrary invariant
completeness. If a candidate polynomial has already been selected, searching
for polynomial multipliers is again linear in those multipliers. Discovering
both sides at once remains the harder problem already identified by Forge.

\subsection{Worked discoveries and proof economy}
For the scaling example the run returned basis $(xy-z)$, transition matrix
$(6)$, and target coefficient $-1$. That complete certificate occupies 121
bytes under the experiment's compact JSON convention. The different sign from
the motivating polynomial is immaterial: the checker verifies the actual basis
and target combination rather than relying on an expected spelling.

For $T(x,y)=(y,2x)$, the basis was $(y,x)$ and the matrix was
$\bigl(\begin{smallmatrix}0&2\\1&0\end{smallmatrix}\bigr)$. The selected target
$x$ uses the coordinate vector $(0,1)$. This demonstrates why a family-valued
certificate is more natural than independent scalar certificates.

For
\[
 (x,y)_0=(0,0),\qquad T(x,y)=(x+1,y+2x+1),
\]
the six-dimensional quadratic feature space yields dimensions
\[
 5\longrightarrow4\longrightarrow3\longrightarrow2\longrightarrow1
 \longrightarrow1.
\]
The surviving basis is $(x^2-y)$. This example is intentionally easy to prove
once the relation is known; it tests the discovery process and the stabilization
criterion. Unlike fitting a polynomial through a few observed states, the final
step checks closure on every state symbolically.

\subsection{Guards, multiple locations, and implementation limits}
Ignoring guards preserves soundness because it adds behaviors. It may lose
nearly every interesting invariant. The experiment includes an initially zero
counter with transition $x\mapsto x+1$ and an intended guard $x<0$. Forgetting
the guard makes $x=0$ fail, although the guarded system cannot move. The worker
returns \unknown{}; it does not refute the guarded source theorem.

A multi-location control-flow graph admits a direct extension: keep a space
$V_\ell$ at each location, intersect an edge source space with the inverse
image of its destination space, and impose initialization only at entry
locations. The sum of the location-space dimensions decreases at every strict
update, giving termination. A certificate has a basis per location and a
closure matrix per edge. This extension is specified here but not implemented
in the delivered single-location prototype.

The implementation caps the monomial feature count at 256. There are at most
$N$ strict descent steps, each involving exact elimination on matrices built
from $N$ features and $M$ ambient monomials. Polynomial-time rational-operation
counts do not imply small bit complexity: coefficients can grow sharply.
For production, exploit sparse supports, cache branch pullbacks, and substitute
an exact matrix backend such as SymPy's \code{DomainMatrix} with rational or
fraction-free row reduction~\cite{domainmatrix}. A target-guided feature set can
be cheaper, but its weaker completeness claim must travel with it.

\section{Hypergeometric antidifference certificates}
\label{sec:gosper}
\status{A denominator-enumerating, coefficient-linear search and an independent
rational-identity checker are implemented. Fifteen positive cases and two
expected-unknown cases were run. A separate SymPy Gosper adapter passed four
regression inputs. A generic Lean sequence-to-ratio bridge is not implemented.}

\subsection{From polynomial increments to multiplicative structure}
The polynomial finite-difference worker already in Forge is appropriate when
$F(n+1)-F(n)$ is polynomial. It does not, by itself, discover the useful
structure in factorials, geometric terms, or rational products. For example,
\[
 \sum_{k=a}^{b-1} k\,k!=b!-a!\qquad(1\le a\le b)
\]
is a one-line proof after recognizing an antidifference, but its summand is not
a polynomial in $k$. The new worker searches for an antidifference expressed
as a rational multiple of the summand.

A hypergeometric term has a rational consecutive ratio. For proof purposes it
is better to encode that ratio without dividing by the term itself:
\begin{equation}
 B(k)t(k+1)=A(k)t(k),\qquad A,B\in\Q[k].
 \label{eq:ratio}
\end{equation}
The source layer must prove this recurrence on the indices where it will be
used. A zero value of $t(k)$ is not automatically a problem; an unjustified
cancellation of $t(k)$ would be.

Seek $G(k)=R(k)t(k)$ with $G(k+1)-G(k)=t(k)$, writing $R=U/V$.
Clearing denominators gives the exact polynomial identity
\begin{equation}
 A(k)U(k+1)V(k)-B(k)U(k)V(k+1)
   =B(k)V(k)V(k+1).
 \label{eq:gospercert}
\end{equation}
For a fixed candidate denominator $V$ and a degree bound on $U$, the coefficients
of $U$ enter linearly. This is the first useful implementation: exact algebra,
a small certificate, and no floating-point recovery.

\subsection{Search and fixed-ansatz completeness}
For every admitted nonzero $V$, enumerate numerator degrees $0,\ldots,d$.
Substitute $U(k)=\sum_{i=0}^{d}u_i k^i$ into
\eqref{eq:gospercert}, equate coefficients, and solve the rational linear
system. A solution yields a proposed pair $(U,V)$.

\par\noindent\begin{minipage}{\linewidth}
\begin{lstlisting}
for V in nonzero_denominator_candidates:
    for numerator_degree in 0 .. degree_cap:
        U := polynomial_with_unknown_coefficients(numerator_degree)
        residual := A * shift(U) * V - B * U * shift(V)
                    - B * V * shift(V)
        coefficients := exact_linear_solve(coefficients(residual) = 0)
        if a solution exists:
            return U[coefficients], V
return UNKNOWN
\end{lstlisting}
\end{minipage}\par\smallskip

Exact elimination is complete for each fixed finite numerator space and fixed
$V$. Enumerating a handful of denominators is not a complete implementation of
Gosper's algorithm. In the delivered bounded search, unsuccessful enumeration
returns \unknown{} even when the mathematical nonexistence of an
antidifference may be known by some stronger method.

Production should use the established algorithm rather than turn the
hand-selected denominator list into a permanent limitation. SymPy exposes
\code{gosper\_normal}, \code{gosper\_term}, and \code{gosper\_sum}; these are
concrete, available search interfaces~\cite{sympy-concrete}. The delivered
\code{sympy\_gosper\_oracle} obtains the consecutive ratio with
\code{hypersimp}, obtains the proposed multiplier with
\code{gosper\_term}, and converts the numerator and denominator to the same
certificate format. It refuses non-polynomial residual components. The
independent identity checker, not the API's naming convention, decides whether
the returned object has the required multiplier meaning.

On the locally tested SymPy 1.14.0, the outputs for
$k k!$, $1/(k(k+1))$, $2^k$, and $k2^k$ were respectively
$1/k$, $-k-1$, $1$, and $(k-2)/k$. Each passed
\eqref{eq:gospercert}. These four adapter checks are reported separately from
the 65 stored corpus certificates; they are not four additional stored results.

\subsection{The theorem that connects the identity to a sum}
\begin{theorem}[Antidifference lifting]
\label{thm:antidifference}
Let $t$ take values in a characteristic-zero field $K$. At a particular index
$k$, suppose \eqref{eq:ratio} holds and $B(k)$, $V(k)$, and $V(k+1)$ are
nonzero. If \eqref{eq:gospercert} is a polynomial identity, then
\[
 \frac{U(k+1)}{V(k+1)}t(k+1)
 -\frac{U(k)}{V(k)}t(k)=t(k).
\]
\end{theorem}
\begin{proof}
Multiply the proposed equality by the nonzero product
$B(k)V(k)V(k+1)$. The left side becomes
\[
 B(k)U(k+1)V(k)t(k+1)-B(k)U(k)V(k+1)t(k).
\]
Use \eqref{eq:ratio} in the first term, factor out $t(k)$, and apply
\eqref{eq:gospercert}. The result is the right side multiplied by the same
nonzero product. Cancel that product. No cancellation of $t(k)$ was used.
\end{proof}

\begin{corollary}[Finite telescoping]
If the hypotheses of Theorem~\ref{thm:antidifference} hold for every integer
$k$ with $a\le k<b$, then
\[
 \sum_{k=a}^{b-1}t(k)=G(b)-G(a),\qquad
 G(k)=\frac{U(k)}{V(k)}t(k).
\]
The empty range $a=b$ has sum zero.
\end{corollary}
\begin{proof}
Replace each summand by $G(k+1)-G(k)$ and cancel adjacent terms, or induct on
the number of summands. The endpoint values are defined by the same justified
formula or separately proved extensions.
\end{proof}

The Python checker validates only \eqref{eq:gospercert} and its syntactic
well-formedness. It does not establish the source recurrence, the quantified
nonvanishing conditions, or the source indexing convention. Those are inputs
to the lifting theorem, not consequences of a successful algebraic calculation.

In Lean, natural-number subtraction is truncated and field division is total.
Thus translating a denominator or a shifted index by surface syntax alone is
unsafe. Prove the integer or rational cast identities, obtain sign or
nonvanishing facts using the source bounds, and only then invoke the
antidifference theorem. The finite range must also be explicit: Mathlib's
\code{Finset.range N} excludes $N$, whereas SymPy's displayed finite-sum bounds
are inclusive and its generalized reversed-bound conventions are not Lean
finite-set semantics~\cite{finset-range,sympy-concrete}.

\subsection{Three representative discoveries}
For $t(k)=k k!$, the recurrence has $A=(k+1)^2$ and $B=k$. The worker tries
constant-denominator numerators through its bound, then finds
\[
 U=1,\quad V=k,\quad R=1/k,\quad G(k)=k!\quad(k\ge1).
\]
The identity check succeeds with a 65-byte compact certificate. Extending this
formula blindly to $k=0$ would replace an intended factorial antidifference by
an undefined rational expression before Lean totalizes it. Either work on
$k\ge1$ or first prove the cancellation and use the nonsingular expression
$G(k)=k!$ on the larger domain.

For $t(k)=1/(k(k+1))$, use $A=k$, $B=k+2$. The worker finds $R=-(k+1)$,
so $G(k)=-1/k$ on $k\ge1$. The familiar reciprocal telescoping formula is
therefore recovered through the same certificate language as the factorial
example, not by a special-purpose syntactic rewrite.

For $t(k)=r^k$ with fixed nonzero rational $r\ne1$, the result is
$R=1/(r-1)$. For $t(k)=kr^k$, a valid multiplier is
\[
 R(k)=\frac{1}{r-1}-\frac{r}{k(r-1)^2}.
\]
At $k=0$ this multiplier is singular, although the product can often be
rewritten to a nonsingular $G$. Treat that cancellation as a separate algebraic
lemma rather than declaring the guard unnecessary because a numerical test
happened to pass.

The corpus includes six rational choices of $r$ for each family, including a
negative base and $r=1/2$. It also keeps bounded-search failures for the
factorial and harmonic summands. These are expected \unknown{} outcomes,
not theorem refutations.

\section{Creative telescoping with a checked boundary plan}
\label{sec:creative}
\status{An order-one, six-unknown coefficient search is implemented and returns
an exact polynomial certificate. The nonsingular boundary argument is derived
mathematically below and checked on 527 concrete index pairs. No general Lean
creative-telescoping adapter or universal binomial-sum proof was run.}

\subsection{Search for a recurrence, not a closed form directly}
A sum may have no convenient hypergeometric antidifference in its summation
variable, while its dependence on a parameter satisfies a short recurrence.
For a bivariate term $t(n,k)$, seek
\begin{equation}
 c_0(n)t(n,k)+c_1(n)t(n+1,k)
   =G(n,k+1)-G(n,k),\qquad G(n,k)=R(n,k)t(n,k).
 \label{eq:creative}
\end{equation}
Summing can then remove the right side. More generally one can admit several
shifts in $n$; the delivered prototype uses order one only.

Suppose the two rational ratios are $A_n/B_n$ and $A_k/B_k$, and write
$R=U/V$. The independently checked polynomial identity is
\begin{multline}
 (c_0B_n+c_1A_n)B_k V V^{+k}\\
 =B_n\bigl(A_k U^{+k}V-B_k U V^{+k}\bigr),
 \label{eq:creativepoly}
\end{multline}
where the superscript $+k$ shifts only $k$. The denominator template $V$ is
fixed during search, and unknown coefficients in $c_0,c_1,U$ enter linearly.
The checker explicitly rejects recurrence coefficients containing $k$; such
coefficients could not simply be taken outside the sum.

This is a concrete entry point for a later Ore-algebra worker. Sage's optional
\code{ore\_algebra} package documents creative telescoping, rational-solution
methods, and desingularization~\cite{ore-algebra}. It is a search dependency,
not a Lean theorem library. A returned recurrence alone is insufficient: an
adapter must obtain a telescoping certificate or another independently
replayable derivation, together with a domain and boundary argument.

\subsection{Discovering the squared-binomial recurrence}
Take
\[
 t(n,k)=\binom nk^2,\qquad S(n)=\sum_{k=0}^{n}\binom nk^2.
\]
On the nonsingular interior the ratios are
\[
 \frac{t(n+1,k)}{t(n,k)}=\frac{(n+1)^2}{(n-k+1)^2},\qquad
 \frac{t(n,k+1)}{t(n,k)}=\frac{(n-k)^2}{(k+1)^2}.
\]
Use the ansatz
\[
 c_0=an+b,\quad c_1=n+c,\quad
 U=k^2(dn+ek+f),\quad V=(n-k+1)^2.
\]
Fixing the coefficient of $n$ in $c_1$ to one removes the irrelevant overall
scaling for this selected template. It does not make the search complete over
all telescopers.

Equating coefficients in \eqref{eq:creativepoly} produced 48 equations in six
unknowns, of rank six. The discovered solution was
\begin{equation}
 c_0=-4n-2,\qquad c_1=n+1,\qquad
 R(n,k)=\frac{k^2(2k-3n-3)}{(n-k+1)^2}.
 \label{eq:binomcert}
\end{equation}
The compact certificate is 313 bytes. The checker reproduces the complete
bivariate polynomial identity by its own sparse arithmetic.

This is an executed search: the coefficients in \eqref{eq:binomcert} were not
inserted as a ready-made certificate. The denominator and low-degree template
were selected in advance. That template selection is an honest form of domain
knowledge, and its cost and limitations should not disappear in a future
benchmark.

\subsection{Why algebraic replay is not yet a sum proof}
The displayed $R$ has a pole at $k=n+1$. In addition, the shifted denominator
$V(n,k+1)$ vanishes at $k=n$. A rational identity derived where denominators are
nonzero cannot be evaluated across these points merely by clearing denominators
and then forgetting the guards.

The summand is zero outside its combinatorial support, but ``zero times a
singular rational expression'' is not a justified boundary value. In Lean the
division expression has a totalized value; that value need not be the extension
required by the telescoping proof. The regression suite includes a test where
the algebraic certificate is accepted while its denominator evaluates exactly
to zero at the boundary. This is intentional: the algebraic checker is doing
its narrower job, not lying about a boundary theorem.

\subsection{A nonsingular flux}
The ratio between consecutive rows of Pascal's triangle suggests rewriting
$R t$ before evaluating it at the singular points. Define, for natural $n,k$,
\begin{equation}
 \widetilde G(n,k)=
 \frac{k^2(2k-3n-3)}{(n+1)^2}\binom{n+1}{k}^2,
 \label{eq:nonsingularflux}
\end{equation}
interpreting the polynomial arithmetic in $\Q$ or $\R$, and setting the
binomial coefficient to zero outside its natural support. The only denominator
is $(n+1)^2$, which is nonzero for natural $n$.

On $0\le k<n$, both adjacent flux values agree with the rational certificate
where the relevant denominators are nonzero, so
\begin{equation}
 (n+1)t(n+1,k)-2(2n+1)t(n,k)
 =\widetilde G(n,k+1)-\widetilde G(n,k).
 \label{eq:fluxlaw}
\end{equation}
The remaining cases $k=n$ and $k=n+1$ must be proved directly.

At $k=n$,
\[
 \widetilde G(n,n)=-n^2(n+3),\qquad
 \widetilde G(n,n+1)=-(n+1).
\]
The left side of \eqref{eq:fluxlaw} is
$(n+1)^3-2(2n+1)$; its right side is $n^2(n+3)-(n+1)$. They agree by polynomial
expansion, including $n=0$.
At $k=n+1$, the left side is $n+1$ because $t(n,n+1)=0$ and
$t(n+1,n+1)=1$. The right side is also $n+1$, since
$\widetilde G(n,n+2)=0$ and $\widetilde G(n,n+1)=-(n+1)$.
Finally $\widetilde G(n,0)=0$.

\begin{theorem}[Squared-binomial recurrence]
For every natural $n$,
\begin{equation}
 (n+1)S(n+1)=2(2n+1)S(n),\qquad S(0)=1.
 \label{eq:binomrec}
\end{equation}
Consequently $S(n)=\binom{2n}{n}$.
\end{theorem}
\begin{proof}
Equation~\eqref{eq:fluxlaw} holds for $k=0,\ldots,n+1$ by the interior
identity and the two direct boundary checks. Sum over precisely that range.
The flux telescopes to
$\widetilde G(n,n+2)-\widetilde G(n,0)=0$. The term $t(n,n+1)$ vanishes, so
the two sums on the left are $(n+1)S(n+1)$ and $2(2n+1)S(n)$.
The initial value is immediate.

The central binomial coefficient $C(n)=\binom{2n}{n}$ has $C(0)=1$ and,
from its factorial expression,
\[
 \frac{C(n+1)}{C(n)}
 =\frac{(2n+2)(2n+1)}{(n+1)^2}
 =\frac{2(2n+1)}{n+1}.
\]
Induction using the nonzero coefficient $n+1$ gives $S(n)=C(n)$.
\end{proof}

The theorem is an ordinary mathematical derivation, not a claim that the Python
boundary loop proved it. The 527 pointwise checks for $0\le n\le30$ and
$0\le k\le n+1$, together with 31 sum and 31 flux checks, are independent
regression evidence for the formulas and the endpoint indexing.

\subsection{A production boundary planner}
A usable summation worker should return a plan with four parts: the algebraic
identity; a region where the source ratio equations and all denominator guards
are proved; a finite list of excluded boundary strata with direct proofs; and
a final range-combination proof. The controller may schedule these as separate
obligations, but none may be omitted merely because the main identity passed.

The first practical domain is rational and factorial/binomial terms on a single
finite natural interval. Factor denominators, solve their integer zero
conditions where possible, and split off the finitely many exceptional indices.
Then attempt cancellation, a neighboring hypergeometric representation, or a
direct recurrence proof for those indices. The squared-binomial example uses
the neighboring-row representation \eqref{eq:nonsingularflux}.

An apparent singularity depending on an unconstrained parameter is not always
a finite boundary issue; it may require parameter case splits. A leading
recurrence coefficient can also vanish and break uniqueness. The worker must
return those cases to the proof planner rather than divide through them.
Infinite sums require convergence and limit arguments and are outside this
proposal's first implementation.

\section{Two-sided certificates for noncommutative algebra}
\label{sec:nc}
\status{Bounded exact search and an independent word-polynomial checker are
implemented. Twenty-five positive instances and two expected-unknown cases
were executed. Two integral certificates were emitted as ordinary Lean replay
scripts; they were not compiled.}

\subsection{Relations are more than normal forms}
Commutative polynomial reasoning identifies $xy$ and $yx$ at the representation
level. That is invalid for matrices, linear operators, and a general associative
algebra. An oriented rewrite system avoids that error, but it introduces a
different burden: useful relations may not have a convenient terminating
orientation, and the required consequence may use a relation in several left
and right contexts.

Mathlib already has \code{noncomm\_ring}. At Forge's exact Mathlib pin, its
source expands and reassociates expressions using simplification and additive
normalization; it supports supplied rewrite lemmas and warns that numeral
powers are fully unfolded. It is an appropriate replay tool for small
identities, not a reason to claim that noncommutative normalization is missing
from Lean~\cite{noncomm-ring}. The increment here is systematic \emph{relation
consequence discovery} with explicit two-sided witnesses.

\subsection{Free associative representation}
Let
\[
 F=\Q\langle X_1,\ldots,X_s\rangle
\]
be the free associative unital algebra. A monomial is a finite word, including
the empty word $\varepsilon$ for $1$. Multiplication concatenates words; it
never sorts letters. A polynomial is a finite rational linear combination of
words.

The source problem consists of relation polynomials $r_1,\ldots,r_m$ asserted
to be zero and a target polynomial $f$ to prove zero. A certificate is a finite
list of terms
\begin{equation}
 f=\sum_{\ell=1}^q c_\ell u_\ell r_{j_\ell}v_\ell,
 \qquad c_\ell\in\Q,
 \label{eq:nccert}
\end{equation}
where $u_\ell,v_\ell$ are words. The checker expands the right side in the free
algebra and compares it with the independently supplied target.

\begin{theorem}[Two-sided certificate soundness]
\label{thm:nc}
Let $A$ be an associative unital $\Q$-algebra, and let an assignment of the
letters to $A$ make every $r_j$ vanish. If \eqref{eq:nccert} holds in the free
algebra, the assignment also makes $f$ vanish.
\end{theorem}
\begin{proof}
Evaluate words by ordered multiplication and extend linearly using the central
rational scalars. This is the canonical algebra homomorphism from the free
algebra. Each evaluated summand on the right is a scalar multiple of
$u_\ell\,0\,v_\ell$, hence zero. Applying the homomorphism to the identity
proves the result.
\end{proof}

Mathlib's \code{FreeAlgebra} supplies a natural semantic setting for a later
reflected checker~\cite{free-algebra}. For a small first implementation, the
same theorem can be instantiated without constructing a generic quotient
algebra: prove the displayed identity with \code{noncomm\_ring}, then rewrite
each relation occurrence to zero.

\subsection{Exact bounded search}
Choose a maximum context degree $D$ and enumerate
\begin{equation}
 \mathcal C_D=
 \{u r_j v: |u|+\deg(r_j)+|v|\le D\}.
 \label{eq:ncdictionary}
\end{equation}
Each member is a column in the rational coefficient space of words of length
at most $D$. Solve the linear membership problem
\[
 A_D c=f.
\]
The implementation inserts columns incrementally into a sparse echelon basis
and records every reduced column's linear combination of original dictionary
columns. After an independent column is inserted, it tries to reduce the target.
When the residual becomes zero, the recorded combination is already the
certificate \eqref{eq:nccert}.

\par\noindent\begin{minipage}{\linewidth}
\begin{lstlisting}
pivots := empty sparse echelon basis
for degree in 0 .. D:
    for relation r and left/right words of the required lengths:
        column := left * r * right
        representation := the corresponding dictionary basis vector
        reduce column using prior pivots, updating representation
        if nonzero:
            normalize its pivot and store both objects
            reduce target through the current pivots
            if target residual is zero:
                return its combination of original columns
return UNKNOWN
\end{lstlisting}
\end{minipage}\par\smallskip

The zero target is handled separately with an empty certificate, including
when there are no relations. This small case matters: a membership algorithm
that forgets it cannot claim completeness even for its own finite span.

\begin{proposition}[Exact finite-span completeness]
Ignoring column and arithmetic resource caps, the bounded worker finds a
certificate exactly when $f\in\Span_\Q\mathcal C_D$.
\end{proposition}
\begin{proof}
Sparse echelon insertion preserves the span of every prefix of input columns.
After all columns are processed, the stored pivot columns span the whole
dictionary. Exact reduction has zero residual exactly for a vector in that
span. The tracked representations express every pivot, and therefore the
accepted target, in the original columns.
\end{proof}

This statement is deliberately weaker than ``complete for ideal membership of
targets of degree at most $D$.'' A low-degree target can require cancellations
between higher-degree ideal multiples. Membership in a bounded dictionary does
not exhaust such possibilities. If all finite degrees are eventually explored,
a particular finite representation will eventually be included, but the
prototype has explicit caps and no unbounded termination claim.

\subsection{An explicit bound on dictionary growth}
For a relation of degree $d_j$, and an alphabet of size $s$, the number of
word-context pairs with total length $e$ is $(e+1)s^e$: choose the left length,
then choose the $e$ letters. Therefore the unreduced dictionary has at most
\begin{equation}
 C_D=\sum_{j=1}^m\sum_{e=0}^{D-d_j}(e+1)s^e
 \label{eq:nccost}
\end{equation}
columns, with a negative upper limit interpreted as an empty sum. The ambient
word space has dimension at most $\sum_{e=0}^{D}s^e$.
Thus exact linearity is not a polynomial-in-$D$ scalability claim. Sparse
support and early target membership can help enormously, but an unrestricted
letter alphabet or degree increase is expensive.

The prototype caps columns at 15,000 and uses degrees up to nine in the main
corpus. Target word supports can guide a later enumeration, but any pruning
that removes contexts loses the displayed finite-span completeness guarantee
unless its admissibility is proved. A practical controller should request a
small $D$, examine the residual and support growth, and escalate only when that
information justifies the cost.

\subsection{Worked certificate: commuting idempotents}
Let $a^2=a$, $b^2=b$, and $ab=ba$. Define $e=a+b-ab$. We want $e^2=e$.
Put
\[
 r_0=a^2-a,\qquad r_1=b^2-b,\qquad r_2=ab-ba.
\]
The search found the following eight-term certificate:
\begin{align}
 e^2-e={}&r_0+r_1-r_0b-br_0-ar_1-r_2a
             +br_0b+r_2ab.\label{eq:idemcert}
\end{align}
No commutativity was assumed while checking this identity: $r_2$ is an explicit
source relation. The checker expands all words on both sides and compares
coefficients. The run considered 41 columns, retained rank 26, and returned a
475-byte compact certificate.

The exact relation indices and contexts are present in
\path{results/certificates.json}. The accompanying emitter turns this
certificate into a generic-ring Lean script, since all eight coefficients are
integers. The script first proves \eqref{eq:idemcert} by noncommutative
normalization and then replaces $r_0,r_1,r_2$ by zero. The generated file is
still an uncompiled target; source generation alone does not establish that its
syntax and proof script work with the pinned toolchain.

\subsection{Commutators and a small representation change}
For $r=xy-yx-1$, one can prove for every fixed positive numeral $n$
\begin{equation}
 xy^n-y^n x-ny^{n-1}
   =\sum_{i=0}^{n-1}y^i r y^{n-1-i}.
 \label{eq:weyl}
\end{equation}
The search found these identities for $n=2,\ldots,8$. At $n=8$ it examined
1,793 columns, retained rank 968, and returned eight certificate terms. The
identity itself is small even when blind context enumeration is much larger.
This is a useful case for learning a checked commutator template after the first
few successes, but the stored numeral certificates are not a proof of the
universally quantified statement in $n$.

An important compression opportunity follows from the example. The ordinary
identity
\[
 [x,uv]=[x,u]v+u[x,v]
\]
can be proved once and used to propose a linear-size certificate for repeated
products. The fallback word-space search remains valuable for discovering a
new relation consequence, while a known structural lemma can prevent it from
re-enumerating thousands of irrelevant contexts. This is a concrete
representation-sensitive route from Forge's existing lemma machinery to the
new worker, not a proposal to resaturate distributivity globally.

For a fixed rational $q$, the relation $xy=qyx$ gives
$x^2y^2=q^4y^2x^2$. Five rational choices, including $q=1/2$, were tested. A
symbolic $q$ would require a coefficient extension with centrality and domain
semantics explicitly represented. The current prototype does not infer those
conditions or treat a new noncentral letter as a rational parameter.

\subsection{Preventing accidental commutativity}
Without hypotheses, $xy-yx$ does not vanish in a free algebra. A concrete
counterexample is given by the matrix units
\[
 X=\begin{pmatrix}0&1\\0&0\end{pmatrix},\qquad
 Y=\begin{pmatrix}0&0\\1&0\end{pmatrix}:
 \quad XY=\begin{pmatrix}1&0\\0&0\end{pmatrix},\quad
 YX=\begin{pmatrix}0&0\\0&1\end{pmatrix}.
\]
Similarly $xy=0$ does not imply $yx=0$: take
$X=\bigl(\begin{smallmatrix}1&0\\0&0\end{smallmatrix}\bigr)$ and
$Y=\bigl(\begin{smallmatrix}0&0\\1&0\end{smallmatrix}\bigr)$.
The bounded search returns \unknown{} on both corresponding unsupported goals.
The displayed matrices are separate mathematical counterexamples; the
worker's failure alone is not a refutation.

The checker preserves ordered letter tuples throughout. Reversing a word,
changing a source relation, introducing an unknown letter, changing a relation
index to a Boolean, or replacing a rational coefficient by a nearby rational
is rejected when it invalidates the identity. ``Normalize the words by sorting
them'' is never an allowable optimization.

\subsection{Coefficients, torsion, and available libraries}
A rational certificate naturally applies to a $\Q$-algebra. An integral
certificate applies to an arbitrary associative unital ring, using its central
integer scalars. Clearing denominators in a rational certificate may prove
$Nf=0$ in a general ring without proving $f=0$: cancellation can fail in the
presence of torsion. Therefore the emitter accepts only integral certificates
for its generic \code{[Ring R]} examples. Rational replay needs an actual
rational-algebra structure or a separately justified cancellation theorem.

Mathlib's \code{linear\_combination} is documented for commutative ring
expressions; it should not be used as an unexamined noncommutative backend.
Small word identities belong with \code{noncomm\_ring}; a larger reflected
checker belongs with the semantics of \code{FreeAlgebra}~\cite{linear-combination,
noncomm-ring,free-algebra}.

Sage exposes free associative algebras via Singular's letterplace implementation,
including homogeneous and weighted-homogeneous methods~\cite{letterplace}.
This is a candidate search accelerator, not a drop-in solution for the
inhomogeneous relation $xy-yx-1$. A production adapter must select a supported
homogenization or a genuinely inhomogeneous method, retain representation
witnesses throughout reduction, and dehomogenize with a proof-preserving map.
A backend returning only a zero normal form is not enough to feed
\eqref{eq:nccert}; it must export a trace or the two-sided combination.

The broader literature includes letterplace treatments of nongraded ideals
and noncommutative Gr\"obner computations~\cite{nongraded-letterplace}. None of
that functionality was exercised in this run. The delivered bounded linear
search is intentionally independent of Singular and works directly with
inhomogeneous relation polynomials.

\section{A stronger universal-assertion backend for Leant and Djex}
\label{sec:behavior}

\subsection{Use the existing candidate stream}
Leant's inspected behavioral interface checks an exact rendered candidate
against its type and then attempts the user assertion and its negation using
\code{decide} and bounded \code{simp}. Its own documentation distinguishes a
finite observation from a universal law and treats unsuccessful proof attempts
as inconclusive. Djex already exposes an incremental typed stream through the
\code{runDjinnTypedQueryStream} family; its streaming report describes retained
candidate ownership, cumulative budgets, and discovery-order delivery
~\cite{leant-behavior,djex-stream}.

There is no reason to replace either mechanism. Add a stronger proof backend
\emph{after} the exact candidate has been fixed. The request consists of that
candidate, its original type, the assertion with that candidate substituted,
and the same scoped source environment. The backend tries ordinary proof
simplification first, then the new invariant or summation worker when a
certified extraction is available. Failure still means inconclusive, not that
the candidate is wrong.

The point of transfer is specific. A universally quantified assertion about a
polynomial accumulator can require an invariant that simplification alone does
not invent. A universal assertion about a factorial-weighted sum can require an
antidifference. Neither assertion is established by checking a larger finite
list of examples. The worker instead supplies the missing mathematical lemma
and an induction or finite-sum argument.

\subsection{A source-to-state extraction contract}
For a candidate recursion over natural numbers, an extractor should return a
transition specification together with a Lean theorem connecting the actual
candidate to the specified orbit. It is not enough to recognize a familiar
function name or print the candidate's apparent update as a polynomial.

For example, suppose the source candidate computes a pair of accumulators by
\[
 (a,b)\longmapsto(2a,3b),
\]
and a derived output is intended to equal their product. Introduce a third
state coordinate for that output, extract its recurrence, and prove that the
source program's state evolves according to the extracted map. An invariant
such as $z-ab=0$ can then prove a universal specification through
Theorem~\ref{thm:subspace}.

A product-program variant compares an implementation and a reference
computation. Carry both states and the common input position in one state tuple;
initialize them under their actual relation; step them in a certified aligned
fashion; and prove that the desired output difference is in the returned
invariant span. A returned equality between state coordinates is useful only
if the output projection theorems connect those coordinates to the source
values.

This first extractor should be syntactically restrictive: structurally
recursive natural functions or explicit folds, polynomial state updates over
rational or real coordinates, and source equations available as proved lemmas.
Integer and natural implementations require checked embedding and operations
with compatible semantics. General recursion, higher-order continuations,
truncated division, and hidden effects remain opaque unless separately modeled.
The delivered Python prototypes consume already extracted problems; they do not
implement this Lean compiler pass.

\subsection{Uniform inputs: a concrete extension for list folds}
\label{sec:uniform}
A list fold has a new element at each step, not just a choice among finitely
many constant transitions. There is a useful extension of the subspace
algorithm that handles an arbitrary fresh polynomial input without discretizing
it or testing finitely many values.

Let an update be $T(x,e)$, polynomial in state $x$ and input $e$. Expand each
feature pullback in input monomials:
\begin{equation}
 m_i(T(x,e))=\sum_{\beta} (K_\beta m_i)(x)e^\beta.
 \label{eq:inputcoeff}
\end{equation}
The finite collection of coefficient maps $K_\beta:E\to\Q[x]$ is linear in
the feature polynomial. Replace each branch pullback in
\eqref{eq:vnext} by all of its coefficient maps. The same descending-subspace
algorithm applies. Its final certificate is
\begin{equation}
 p(T(x,e))=\sum_\beta H_\beta p(x)e^\beta
          =H(e)p(x).
 \label{eq:parametricmatrix}
\end{equation}
The input is universally quantified afresh at each step, not retained as a
single fixed parameter throughout the execution.

\begin{proposition}[Uniform-input extension]
The coefficient-map construction terminates by dimension descent and proves
its invariant family for every finite sequence of input values. If $E$ contains
all state monomials of degree at most $D$ and $T$ is affine in the state
variables, with polynomial dependence on the input, it is complete for
rational state-polynomial equalities of degree at most $D$ that hold under all
such input sequences and the specified initial states.
\end{proposition}
\begin{proof}
For soundness, each coefficient map sends the final space into itself; hence
\eqref{eq:parametricmatrix} preserves simultaneous vanishing for every input.
Induct on the input sequence length.

For completeness, let $W$ be the degree-bounded polynomials vanishing on all
reachable states. For every reachable $x$, the polynomial
$p(T(x,e))$ is zero for every input $e$. Since the coefficient field embeds in
an infinite field, each coefficient in $e$ vanishes at that $x$. Affineness in
state variables ensures that each coefficient polynomial still has state degree
at most $D$. Thus every coefficient map preserves $W$, and the maximality
argument of Theorem~\ref{thm:subspace} applies.
\end{proof}

This extension is specified and proved here but not implemented in the main
corpus. Its implementation is concrete: extract coefficient matrices in the
fresh input variables, reuse the existing descent routine, and change each
constant transition matrix in the certificate to a finite list of constant
matrices tagged by input monomials. Replay expands the resulting polynomial
matrix. It is not the unrestricted jointly unknown polynomial-multiplier
problem: the input-coefficient maps are fixed by the known transition.

As a small fold example, two accumulators both updated by $e^2$ preserve
$a-b$ uniformly in $e$. If the reference update is instead $2e^2$, the equality
fails except under extra assumptions, which the worker must not invent.
This distinguishes universal accumulator reasoning from finite behavioral
sampling in exactly the setting where Leant's assertion interface is useful.

\subsection{A proposed narrow proof-service interface}
The proposed interface below is new code to implement, not an API claimed to
exist in either neighboring project:
\par\noindent\begin{minipage}{\linewidth}
\begin{lstlisting}
proveCandidateAssertion(exactCandidate, assertion, sourceContext):
    try ordinary Lean proof simplification
    if a source-linked polynomial orbit specification is available:
        request an inductive-subspace certificate
        prove its polynomial identities in Lean
        combine with the orbit theorem and the assertion projection
    if a source-linked finite sum specification is available:
        request a telescoping certificate and a boundary plan
        prove ratios, guards, exceptions, and the finite range equation
    return a checked assertion proof, or Inconclusive(reason)
\end{lstlisting}
\end{minipage}\par\smallskip

A theorem from this backend is attached to the exact candidate already being
considered. It is not a reason to consume the unobserved tail of Djex's stream,
borrow a receipt from a textually identical candidate, or reset a budget after
an inconclusive attempt. Those policies already exist in the inspected code
and should remain unchanged~\cite{leant-behavior,djex-stream}.

\section{Lean integration: concrete modules and library choices}
\label{sec:integration}

\subsection{Implement the replay boundary first}
The first deliverable inside Forge should be an ordinary proof-producing
adapter for one worker, not a modification of the core \grind{} extension
state. Forge already discusses both the conservative and invasive integration
routes; this report relies on that choice rather than repeating the architecture
~\cite{forge-integration}.

A useful new module layout is:
\par\noindent\begin{minipage}{\linewidth}
\begin{lstlisting}
Forge/Invariant/FeatureSpace.lean    -- source-linked feature denotation
Forge/Invariant/ClosureReplay.lean  -- init, matrix step, target identities
Forge/Invariant/ExtractOrbit.lean   -- source recursion -> orbit theorem
Forge/Invariant/SearchBridge.lean  -- bounded rational subspace oracle

Forge/Summation/RatioSpec.lean      -- source sequence ratio theorem
Forge/Summation/Antidifference.lean -- rational identity lifting
Forge/Summation/FiniteRange.lean    -- exact finite sum combinators
Forge/Summation/BoundaryPlan.lean   -- exceptional indices and flux extension
Forge/Summation/SearchBridge.lean  -- Gosper/Ore certificate import

Forge/Noncomm/WordPoly.lean         -- ordered-word polynomial semantics
Forge/Noncomm/TwoSidedReplay.lean   -- sum c * left * relation * right
Forge/Noncomm/SearchBridge.lean     -- bounded word-context discovery
\end{lstlisting}
\end{minipage}\par\smallskip
These paths are proposed additions to Forge; they are not files asserted to
exist in the reviewed repository. The accompanying \path{lean/} directory
contains smaller standalone replay targets, not this full module tree.

\subsection{What each adapter must prove}
For invariants, the source layer supplies an initialization theorem, the
transition interpretation, and the target interpretation. The replayer proves
each component identity in \eqref{eq:initcert}--\eqref{eq:targetcert} with
\code{ring} on the original scalar expressions, constructs simultaneous
vanishing, and applies a reachable-state induction theorem. It need not verify
the producer's row reduction or its claim of maximality to accept a positive
certificate.

For summation, the source layer supplies the actual ratio recurrence and the
finite index set. The replayer checks \eqref{eq:gospercert} or
\eqref{eq:creativepoly}, proves the denominator guards and exceptional cases,
and applies a finite telescoping theorem. The algebraic identity should be
stored as an intermediate theorem even when a boundary obligation remains
unsolved; it is partial progress, not a completed sum proof.

For noncommutative reasoning, the source layer supplies an ordered variable
table and actual relation proofs. A small replayer expands the displayed
combination with \code{noncomm\_ring} and rewrites the exact relation subterms.
A generic reflected implementation can later evaluate a word-polynomial
certificate and apply a soundness theorem over the appropriate algebra.
No commutative solver should be allowed to identify words during this replay.

\subsection{Concrete dependencies and their roles}
\begin{longtable}{@{}p{.24\textwidth}p{.39\textwidth}p{.30\textwidth}@{}}
\toprule
Library or mechanism & Specific use & Status in this work\\
\midrule\endhead
SymPy 1.14.0 & Polynomial coefficient matrices; exact nullspace and linear
solves; tested \code{hypersimp}/\code{gosper\_term} adapter & Installed and
executed\\
Python \code{fractions.Fraction} & Independent sparse commutative and word
polynomial replay & Executed, including a fresh \code{python -S} process\\
SymPy \code{DomainMatrix} & Rational-domain or fraction-free matrix backend for
larger subspace problems & Documented optimization target, not used in timing
results~\cite{domainmatrix}\\
Mathlib \code{ring}, finite sums, \code{noncomm\_ring} & Ordinary proof-term
reconstruction of identities and telescoping & Source-level targets;
\code{noncomm\_ring} inspected at the exact pin~\cite{noncomm-ring}\\
Mathlib \code{FreeAlgebra} & Semantics for a future word-polynomial reflected
checker & Library identified; no compiled checker~\cite{free-algebra}\\
Sage \code{ore\_algebra} & More general telescoper search, rational solutions,
and desingularization & Proposed search backend; not installed or
executed~\cite{ore-algebra}\\
Sage/Singular letterplace & Search acceleration for supported free-algebra
ideal computations with retained witnesses & Proposed; homogeneous and
inhomogeneous modes must not be confused~\cite{letterplace,nongraded-letterplace}\\
\bottomrule
\end{longtable}

Recent Lean work already supplies certificate-oriented commutative polynomial
tactics for Gr\"obner reasoning, ideal membership, and related tasks. It should
be evaluated as a reusable backend when a future invariant worker needs
commutative ideal certificates; it is not evidence that the two-sided
noncommutative adapter above already exists~\cite{groebner-tactics}.
This report does not claim compatibility with that external tactic project or
report its performance as its own.

\subsection{Generated Lean files and the first acceptance gate}
The archive contains a core-only reachable-state induction theorem; ordinary
rational identities and a finite telescoping combinator; and two generated
noncommutative replay scripts. The generator only emits integral certificates
into its generic-ring theorems. These files contain no placeholder axioms or
\code{sorry}, but they were not compiled in this environment. Their current
status remains \code{NOT\_COMPILED}, not ``probably verified.''

Before admitting any worker into Forge, run the real pinned Lake project,
compile these files, inspect the axiom inventories, and exercise the replayer
on negative controls that modify the source problem rather than just the
certificate. Then add a source expression reifier and repeat the controls on
actual Lean goals. A Python certificate checked for a manually prepared
polynomial problem does not test that reifier.

The first public tactic syntax can be explicit rather than magical. For
example, a diagnostic command could show a source-linked transition model,
selected features, and a proposed invariant before the full automatic planner
chooses them. The eventual \forge{} front end can call the same interface
silently. This keeps the first capability test focused on mathematical
coverage instead of conflating it with complete automation of extraction and
scheduling.

\section{Executed experiments and what they establish}
\label{sec:experiments}

\subsection{Environment and method}
The recorded run used Python 3.13.5, SymPy 1.14.0, and seed 20260915 on the
Linux environment identified in \path{results/summary.json}. No Lean or Lake
executable was available. No stock \grind{}, \code{simp}, \code{ring}, or
\code{noncomm\_ring} benchmark was run. Existing repository test results were
not added to these counts.

The experiment script constructs each problem independently of the certificate
returned by search, invokes the checker with both arguments, and saves them in
separate files. After generation, a fresh Python process runs the checker with
\code{-S}, which excludes automatic site-package initialization. The replay
uses only the standard library and imports neither SymPy nor the producer.
This confirms a useful dependency separation; it does not turn Python into a
proof kernel.

Search timing covers the individual search call on an already constructed
symbolic problem, including coefficient-system construction within that call.
It excludes Python/SymPy import startup, corpus generation, mutation tests,
serialization, the conservation ablation, and any Lean work. Replay timing
covers one in-memory certificate check. Each number is one observation, not a
statistical latency estimate.

\begin{table}[htbp]\centering
\small
\begin{tabular}{@{}lrrr@{}}
\toprule
Worker & Certificates & Search (ms) & Replay (ms)\\
\midrule
Inductive subspaces & 24 & 216.1 & 7.59\\
One-variable antidifferences & 15 & 247.0 & 2.19\\
Creative telescoper & 1 & 107.9 & 1.59\\
Two-sided word relations & 25 & 110.5 & 4.24\\
\midrule
Total & 65 & 681.4 & 15.60\\
\bottomrule
\end{tabular}
\caption{Recorded single-run totals for the stored positive corpus. Search and replay are Python operations on prepared problems, not Lean tactic timings. The two summation rows certify algebraic identities only.}
\label{tab:timing}
\end{table}


\subsection{The positive corpus}
The 24 invariant cases consist of six named examples and 18 planted affine
families. The named cases are the scaling graph, the mutually preserved pair,
the square counter, a rotation preserving a circle, two equal accumulators,
and a two-branch scaling graph. The planted families initialize $z=x+y$ and
construct two affine transitions preserving that relation with a scalar factor
different from one. These are mechanism tests chosen to distinguish the two
certificate languages, not a representative sample of arbitrary Lean proofs.

The 15 one-variable summation cases consist of a factorial-weighted term, two
reciprocal terms, six geometric terms, and six weighted geometric terms.
The sole creative-telescoping instance is the squared-binomial example. The
25 noncommutative cases consist of seven Weyl-commutator numeral instances,
five rational $q$-commutation instances, the commuting-idempotent join, and
12 planted two-sided combinations.

All 65 returned certificates passed independent replay. Their acceptance means
exact closure or algebraic identities for the explicitly supplied problem
objects. For the summation kinds it does \emph{not} include a source sequence
bridge or universally checked boundary plan.

\subsection{A same-feature invariant ablation}
The reference ablation uses the same feature degree, initial map, branch maps,
and target, but restricts its candidate space to
\[
 p\circ I=0,\qquad p\circ T_b-p=0\quad\text{for every branch}.
\]
It solves these exact coefficient constraints directly. This is an
implementation of the \emph{conservation-only certificate language}, not a
rerun of an original Forge executable.

On the 24 selected cases the new subspace algorithm proves target membership in
24 cases and the conservation-only ablation in three. The latter succeeds on
the square counter, circle rotation, and equal accumulators. The other examples
were intentionally selected or planted to require scaling or collective
preservation. Thus the observed gap demonstrates added expressive coverage,
not a 24-to-3 estimate for general proof workloads or a performance victory over
Lean's native tactics.

For the scaling graph the gap can be explained mathematically, independently
of the experiment. In degree at most two, the initialization kernel is the
one-dimensional span of $xy-z$. Pullback acts on that span as multiplication
by six. Its only fixed vector is zero, whereas the whole line is invariant.
No better implementation of the same conservation-only ansatz can recover the
missing target in that feature space.

\subsection{Negative and adversarial controls}
Seven unsupported cases are recorded as \unknown{}: three invariant cases,
two bounded antidifference failures, and two noncommutative consequence
failures. These outcomes are never interpreted as global impossibility
certificates. The nonlinear degree-escape case is especially important because
the desired invariant is true; it checks that the report distinguishes
fragment incompleteness from a false goal.

The 188 certificate mutations are assembled from uniform and targeted
controls. Every stored certificate is tested with its essential payload
removed and with an unexpected field inserted. Every invariant certificate is
also tested with a different target and with a missing transition matrix.
Additional cases change a branch, zero a denominator, introduce a floating-point
or noncanonical coefficient, duplicate a monomial, allow a recurrence
coefficient to depend on $k$, change a source relation, misuse a Boolean as an
index, introduce an unknown letter, and make a tiny exact rational change.
Every tested invalid variant is rejected.

These tests do not establish that the decoder is secure against every hostile
byte stream. It checks field sets, rational string canonicality, dimensions,
monomial ordering, bounded indices, degree and size caps, and a product-work
budget. It still relies on Python's parser and allocation behavior, and it is
not fuzzed across all possible nested JSON inputs. The code is a research
checker with explicit limits, not a hardened network service.

\subsection{Independent arithmetic and boundary checks}
One hundred randomized substitution tests compare the checker's sparse
substitution with SymPy's simultaneous substitution and expansion. These
exercise a different algorithmic path from accepting a stored solver result.
They reduce the risk of a shared substitution convention error; they are not
an exhaustive arithmetic proof.

For the squared-binomial flux, 527 concrete $(n,k)$ pairs cover
$0\le n\le30$ and $0\le k\le n+1$. The experiment also checks the finite sum
and the three relevant flux values for every one of those 31 $n$ values.
The $n=0$ case and both singular strata occur in this range. Additional
factorial-sum checks use $1\le a\le8$ and $a\le b\le11$, including empty
ranges. None of these samples is used as the acceptance condition for a
polynomial identity.

The generated run records 1,100 assertions in total; this count includes
individual samples, positive checks, mutation rejections, dimension checks,
and the fresh-process replay. It is \emph{not} a count of 1,100 independent
universal theorems. The separate standard \code{unittest} suite has 13 named
test methods, with four SymPy-adapter inputs inside one method and stored
certificate subtests inside another. Their raw log is included.

\subsection{What has not been tested}
No result here measures automatic Lean term extraction, dependent typing,
proof-term size inside the kernel, equality transport, source cast handling,
the performance of a real Forge controller, or the behavior of a compiled
\grind{} plugin. The uncompiled Lean files do not fill those gaps.
There is also no large benchmark demonstrating that the chosen context or
feature bounds are sensible on typical formalization tasks.

A useful next comparison would separate three conditions: a raw goal, a goal
with the worker's discovered mathematical lemma supplied, and a goal with a
complete certificate supplied. Compare each against the pinned baseline under
identical theorem inventories and trust settings. That comparison can reveal
whether the value lies in discovery, local proof reconstruction, or cross-worker
cooperation. The present experiment establishes the first component's
possibility on explicit problems, not the whole system's superiority.

\section{Implementation sequence and acceptance criteria}
\label{sec:roadmap}

\subsection{First increment: noncommutative relation replay}
This is the smallest isolated production addition. Its semantic objects are
words, rational or integer coefficients, relation references, and a target.
Start with integral certificates in associative rings and small source
expressions. Require compilation of the emitted examples, a checked source
reifier preserving word order, and negative tests where $ab$ and $ba$ differ.
Then connect the existing bounded producer and expose the resulting theorem to
Forge. Add rational coefficients only with an explicit algebra structure.

This ordering is not a prediction that the noncommutative worker has the
largest eventual coverage. It has the shortest path from an already executed
certificate to a useful Lean theorem, with relatively little domain-specific
boundary work.

\subsection{Second increment: source-linked inductive subspaces}
Implement the generic simultaneous-invariant induction theorem and replay
small closure matrices using ordinary polynomial normalization. Support a
single source-linked natural recursion or an explicit transition relation
before general list-fold extraction. Acceptance requires the scaling example,
a mutually preserved family, multiple branches, a genuine target projection,
and failures when the initialization or a branch is changed.

Only after these pass should feature selection become automatic. Add the
uniform-input coefficient-map extension of Section~\ref{sec:uniform} to cover
polynomial list folds without finite-alphabet sampling. Add multi-location
control-flow spaces separately; neither feature requires changing the
positive-certificate theorem into a global completeness claim.

\subsection{Third increment: a small summation library}
Prove the finite-range telescoping combinator, the guarded antidifference
lifting theorem, and ratio lemmas for a deliberately short list of source
constructors. Integrate the tested SymPy Gosper adapter through typed data,
not by importing executable code or an unproved symbolic equality.
Require tests at zero, at empty ranges, at denominator zeros, and at parameter
values where a recurrence coefficient vanishes.

The first complete creative-telescoping acceptance target should be the
squared-binomial theorem with the nonsingular flux, because it exercises the
hard boundary mechanism rather than only an interior rational identity.
General Ore-algebra import comes later, after the adapter can explain exactly
which returned object becomes which Lean theorem.

\subsection{Stopping rules that preserve value}
Each worker should be independently useful before the next is started. The
noncommutative worker can prove a relation consequence without induction. The
invariant worker can export an equality useful to \grind{} or an arithmetic
worker. The summation worker can produce a checked local antidifference even
when a more ambitious closed-form goal needs additional steps.

When a bound is reached, retain any already checked consequences and report
\unknown{} for the remaining request. Do not infer that a theorem is false
because a feature space, denominator list, or context dictionary failed. Do not
increase a bound indefinitely when the failure has a structural explanation,
such as nonlinear degree escape or a singular endpoint. Those explanations are
routing information for Forge's existing planner.

\section{Conclusion}
The strongest next version of Forge need not acquire another general-purpose
search slogan. It can acquire three concrete mathematical capabilities: find
a space of equations that evolves into itself, find a local flux whose finite
sum gives a recurrence, and find a two-sided combination of noncommutative
relations. Their certificates are small, their checking conditions are explicit,
and their failure modes point to meaningful alternative workers.

The delivered prototypes establish that these searches can be executed and
that their proposed identities survive independent exact replay. The invariant
ablation demonstrates a real increase in the chosen certificate language's
coverage. The binomial example demonstrates why source semantics cannot be
collapsed into polynomial checking. The noncommutative examples demonstrate
how to discover useful consequences without falsely assuming commutativity.

The remaining work is substantial but concrete: source-linked Lean extraction,
actual kernel replay, and fair end-to-end evaluation. Those are engineering
acceptance gates, not reasons to replace the algorithms with a vague proposal.

\appendix
\section{Certificate formats and reproducibility}
\label{app:formats}

\subsection{Canonical arithmetic data}
A commutative polynomial is a sorted list of pairs \code{\{m,c\}}. The field
\code{m} is an exponent vector of fixed arity; \code{c} is a canonical rational
string such as \code{"-3/2"}. Zero coefficients and duplicate monomials are
forbidden. A noncommutative polynomial uses the same external shape, but
\code{m} is an ordered word of letter indices. The \code{kind} tag fixes the
interpretation; a word is never decoded as a commutative exponent vector.

An invariant certificate contains the polynomial basis, one rational matrix per
source transition, and coefficients expressing the source target in that
basis. The initial map and all transitions live in the source problem file,
not in the producer's certificate. The scaled-graph certificate has basis
$xy-z$, matrix $[6]$, and target coefficient $-1$.

A Gosper certificate contains only $U,V$; the source problem supplies $A,B$.
The creative certificate adds $c_0,c_1$ and uses the source pair of ratio
numerators and denominators. A two-sided certificate contains terms with
\code{relation}, \code{left}, \code{right}, and \code{coefficient} fields.
The source relation array and target are not replaceable by those terms.

The checker permits no JSON floating-point literals, duplicate object keys,
unknown certificate fields, noncanonical rational strings, or invalid
monomial orderings. It limits source dimensions to eight, polynomial degree
to 32, expanded terms to 20,000, and elementary multiplication-pair work to
two million per certificate. These are conservative operational rejection
limits; changing them does not change the mathematical soundness argument.

\subsection{Archive contents}
\begin{longtable}{@{}p{.36\textwidth}p{.59\textwidth}@{}}
\toprule
Path & Purpose\\\midrule\endhead
\path{article/forge_extensions.tex} and \path{article/forge_extensions.pdf}
& This report and its editable source\\
\path{prototype/search.py} & Untrusted exact search, including the optional
SymPy Gosper adapter\\
\path{prototype/checker.py} & Standard-library-only independent replay\\
\path{prototype/experiments.py} & Deterministic corpus, ablation, mutations,
differential checks, and result capture\\
\path{prototype/test_regressions.py} & Thirteen named regression methods\\
\path{prototype/emit_lean.py} & Integral noncommutative replay-script emitter\\
\path{results/} & Source problems, certificates, per-case measurements,
summary, expected unknowns, assertions, mutations, and raw logs\\
\path{lean/} & Pinned project targets and uncompiled ordinary Lean proofs\\
\path{docs/review-scope.md} and \path{docs/status.json} & Source pins,
nonduplication scope, and machine-readable evidence status\\
\bottomrule
\end{longtable}

\subsection{Commands}
From the archive root, replay the stored certificates without third-party
packages:
\par\noindent\begin{minipage}{\linewidth}
\begin{lstlisting}
python -S prototype/checker.py
\end{lstlisting}
\end{minipage}\par\smallskip
Install the search dependency and run the named tests:
\par\noindent\begin{minipage}{\linewidth}
\begin{lstlisting}
python -m pip install -r prototype/requirements.txt
python -m unittest discover -s prototype -p 'test_*.py' -v
\end{lstlisting}
\end{minipage}\par\smallskip
Regenerate the experiment into a fresh directory:
\par\noindent\begin{minipage}{\linewidth}
\begin{lstlisting}
python prototype/experiments.py --output reproduced-results
python -S prototype/checker.py reproduced-results/problems.json \
    reproduced-results/certificates.json
\end{lstlisting}
\end{minipage}\par\smallskip
The generator refuses a nonempty output directory, so the recorded results are
not silently overwritten. A new run can have different timings while producing
the same exact certificates at the recorded dependency version.
Build this report with two pdf\LaTeX{} passes or with \code{latexmk}:
\par\noindent\begin{minipage}{\linewidth}
\begin{lstlisting}
cd article
latexmk -pdf -interaction=nonstopmode -halt-on-error forge_extensions.tex
\end{lstlisting}
\end{minipage}\par\smallskip
The separate \path{lean/README.md} describes the unexecuted Lean build and
axiom-inspection steps. No \code{lake-manifest.json} is fabricated for a build
that did not occur.

\section{A compact set of regression targets for a Lean port}
\label{app:targets}
A port should preserve the following distinctions before it is called an
implementation of these workers.

\textbf{Invariant positive target.} Starting at $(u,v,uv)$ and applying
$(x,y,z)\mapsto(2x,3y,6z)$ any number of times preserves $z=xy$.
The initial parameters are universally quantified. Changing the third update
to $6z+1$ must invalidate the same certificate.

\textbf{Collective closure target.} Starting at $(0,0)$ and applying
$(x,y)\mapsto(y,2x)$ preserves both coordinates' vanishing. A replay that checks
only one diagonal entry of the matrix must not pass.

\textbf{Nonlinear incompleteness target.} Starting at zero under $x\mapsto x^2$
preserves $x=0$, but the finite constant-closure subspace search returns
\unknown{}. The tactic may use another proof route; it must not call the
theorem false.

\textbf{Antidifference guard target.} The factorial-weight multiplier $1/k$
is valid on the intended positive indices. A proof that substitutes $k=0$
without cancellation or a separate extension theorem is invalid.

\textbf{Creative boundary target.} The squared-binomial certificate must prove
the two exceptional strata $k=n$ and $k=n+1$, not just the generic rational
identity. The $n=0$ case has no generic interior and must still be covered.

\textbf{Word-order target.} Given $ab-ba=1$, prove the fixed numeral
commutator consequence. With no relation, do not prove $ab=ba$. Given only
$ab=0$, do not prove $ba=0$.

\textbf{Coefficient-domain target.} Replay an integral two-sided certificate
in an arbitrary associative ring. Require the appropriate rational scalar
structure before replaying a genuinely rational one; do not cancel a cleared
integer denominator in a torsion ring.

\textbf{Exact-candidate target.} A universal behavioral proof must concern the
same source term that was type-checked by Leant. Failure of the new backend is
inconclusive and does not consume a success slot or establish noninhabitation.

\begin{thebibliography}{99}
\small
\bibitem{forge} Vladimir Reshetnikov. \emph{Forge}, README and repository tree.
Revision \nolinkurl{674521027d968d59f7b83220ed52304a85cb55e2}, inspected September
15, 2026. \url{https://github.com/VladimirReshetnikov/Forge/tree/674521027d968d59f7b83220ed52304a85cb55e2}.
\bibitem{forge-structural} Forge, \emph{Structural search section}. \path{article/sections/05-structural.tex}, same revision
as~\cite{forge}.
\bibitem{forge-nonlinear} Forge, \emph{Nonlinear arithmetic section}. \path{article/sections/06-nonlinear.tex}, same
revision as~\cite{forge}.
\bibitem{forge-external} Forge, \emph{Boolean search, equality exploration, and
external engines}. \path{article/sections/07-boolean-and-external.tex}, same
revision as~\cite{forge}.
\bibitem{forge-integration} Forge, \emph{Lean integration: a conservative bridge
before an invasive one}. \path{article/sections/08-integration.tex}, same
revision as~\cite{forge}.
\bibitem{forge-provenance} Forge, \emph{Provenance: what came from where}.
\path{article/sections/B-provenance.tex}, and
\path{docs/IDEAS-FROM-LEANT-DJEX.md}, same revision as~\cite{forge}.
\bibitem{leant-behavior} Vladimir Reshetnikov. Leant, \emph{Behavioral assertions
during synthesis}. \path{docs/behavioral-synthesis.md}, revision
\nolinkurl{6bf05ad78c467989e68290f2d08bbed40802d485}.
\url{https://github.com/VladimirReshetnikov/Leant/blob/6bf05ad78c467989e68290f2d08bbed40802d485/docs/behavioral-synthesis.md}.
\bibitem{djex-stream} Vladimir Reshetnikov. Djex, \emph{Incremental Djinn
behavioral synthesis}. Report dated September 6, 2026, revision
\nolinkurl{e8778f4ebd63e1f9b9b410fa4de8d14a8a04c9e5}.
\url{https://github.com/VladimirReshetnikov/Djex/blob/e8778f4ebd63e1f9b9b410fa4de8d14a8a04c9e5/docs/reports/2026-09-06-djinn-behavioral-streaming.md}.
\bibitem{affine-invariants} Ehud Hrushovski, Jo\"el Ouaknine, Amaury Pouly, and
James Worrell. \emph{Polynomial invariants for affine programs}. LICS 2018,
pp.~530--539. DOI: 10.1145/3209108.3209142.
\url{https://arxiv.org/abs/1802.01810}.
\bibitem{domainmatrix} SymPy Development Team. \emph{Introducing the
DomainMatrix of the poly module}. SymPy 1.14.0 documentation.
\url{https://docs.sympy.org/latest/modules/polys/domainmatrix.html}.
\bibitem{sympy-concrete} SymPy Development Team. \emph{Concrete mathematics},
including \code{hypersimp}, \code{gosper\_normal}, \code{gosper\_term}, and
\code{gosper\_sum}. SymPy 1.14.0 documentation.
\url{https://docs.sympy.org/latest/modules/concrete.html}.
\bibitem{finset-range} Mathlib contributors. \emph{Mathlib.Data.Finset.Range}.
\url{https://leanprover-community.github.io/mathlib4_docs/Mathlib/Data/Finset/Range.html}.
\bibitem{ore-algebra} Manuel Kauers and collaborators. \emph{ore\_algebra}:
Sage implementation of Ore algebras, Ore polynomials, and differentially finite
functions. Official Sage package documentation and project.
\url{https://doc.sagemath.org/html/en/reference/spkg/ore_algebra.html};
\url{https://github.com/mkauers/ore_algebra}.
\bibitem{noncomm-ring} Jireh Loreaux, Kim Morrison, Oliver Nash, and Mathlib
contributors. \emph{Mathlib.Tactic.NoncommRing}. Source inspected at
\nolinkurl{1cf325a0cf67aca2b04d76b5380ff6a9e410aefa}.
\url{https://github.com/leanprover-community/mathlib4/blob/1cf325a0cf67aca2b04d76b5380ff6a9e410aefa/Mathlib/Tactic/NoncommRing.lean}.
\bibitem{linear-combination} Mathlib contributors.
\emph{Mathlib.Tactic.LinearCombination}. Official generated documentation.
\url{https://leanprover-community.github.io/mathlib4_docs/Mathlib/Tactic/LinearCombination.html}.
\bibitem{free-algebra} Mathlib contributors. \emph{Mathlib.Algebra.FreeAlgebra}.
Official generated documentation.
\url{https://leanprover-community.github.io/mathlib4_docs/Mathlib/Algebra/FreeAlgebra.html}.
\bibitem{letterplace} SageMath contributors, including Simon King.
\emph{Free associative unital algebras, implemented via Singular's letterplace
rings}, and associated element documentation. Official Sage reference.
\url{https://doc.sagemath.org/html/en/reference/algebras/sage/algebras/letterplace/free_algebra_letterplace.html}.
\bibitem{nongraded-letterplace} Roberto La Scala. \emph{Extended letterplace
correspondence for nongraded noncommutative ideals and related algorithms}.
2012. \url{https://arxiv.org/abs/1206.6027}.
\bibitem{groebner-tactics} Hao Shen, Junyu Guo, Junqi Liu, and Lihong Zhi.
\emph{Automated Tactics for Polynomial Reasoning in Lean 4}. April 15, 2026,
arXiv:2604.13514v1. \url{https://arxiv.org/html/2604.13514v1}.
\end{thebibliography}
\end{document}
